\documentclass[preview]{standalone}

\usepackage{amsmath}
\usepackage{amssymb}
\usepackage{stellar}
\usepackage{definitions}

\begin{document}

\id{psicologia-giudizio-decisione}
\genpage

\section{Giudizio e decisione}

\begin{snippet}{giudizio-decisione-introduzione}
    Giudizio e presa di decisione applicano risorse cognitive limitate a
    situazioni incerte. Il giudizio riguarda la formazione di valutazioni,
    mentre la decisione riguarda la scelta comportamentale tra opzioni
    disponibili.
\end{snippet}

\begin{snippetdefinition}{giudizio-definition}{Giudizio}
    Il \emph{giudizio} è il processo attraverso cui si formano opinioni, si
    raggiungono conclusioni e si fanno valutazioni critiche di eventi e persone.
\end{snippetdefinition}

\begin{snippetdefinition}{presa-decisione-definition}{Presa di decisione}
    La \emph{presa di decisione} è il processo con cui si sceglie tra due o più
    alternative, accettando o rifiutando le opzioni disponibili.
\end{snippetdefinition}

\begin{snippetdefinition}{euristica-giudizio-definition}{Euristica}
    Un'\emph{euristica} è una regola empirica informale, cioè una scorciatoia
    cognitiva che riduce la complessità nella formulazione del giudizio.
\end{snippetdefinition}

\begin{snippet}{cassetta-attrezzi-adattativa}
    Le euristiche possono essere interpretate come una cassetta degli attrezzi
    adattativa: un repertorio di procedure rapide ed economiche che spesso
    portano a giudizi efficaci, ma che in alcune condizioni producono errori
    sistematici.
\end{snippet}

\begin{snippetdefinition}{euristica-disponibilita-definition}{Euristica della disponibilità}
    L'\emph{euristica della disponibilità} è una scorciatoia di giudizio in cui
    la valutazione dipende dalla facilità con cui le informazioni vengono
    recuperate dalla memoria e dai contenuti che sembrano più disponibili.
\end{snippetdefinition}

\begin{snippetdefinition}{euristica-rappresentativita-definition}{Euristica della rappresentatività}
    L'\emph{euristica della rappresentatività} è una scorciatoia di giudizio per
    cui si assume che qualcosa appartenga a una categoria se possiede le
    caratteristiche tipiche dei membri di quella categoria.
\end{snippetdefinition}

\begin{snippetexample}{avvocato-sport-example}{Avvocato e sport preferito}
    Se una persona viene descritta come un avvocato di successo, competitivo,
    molto alto, controllato e vanitoso, si può essere portati a scegliere il
    tennis come suo sport preferito. Tuttavia, una risposta più probabile è
    ``un gioco con la palla'', perché include anche il tennis. L'errore nasce
    dall'attrattiva di un'alternativa molto rappresentativa.
\end{snippetexample}

\begin{snippetexample}{acqua-gelida-example}{Esperimento sull'acqua gelida}
    In un esperimento, i partecipanti affrontavano due prove dolorose. Nella
    prova breve tenevano la mano per \(60\) secondi in acqua a \(14^\circ\)C.
    Nella prova lunga, dopo i primi \(60\) secondi, tenevano la mano per altri
    \(30\) secondi mentre l'acqua raggiungeva \(15^\circ\)C. Quando veniva
    chiesto quale prova preferissero ripetere, molti sceglievano quella lunga,
    perché il momento finale era meno doloroso.
\end{snippetexample}

\includesnpt[width=55\%,src=/snippet/static-psychology/cold-pressor-peak-end-rule.jpg]{centered-img}

\begin{snippetdefinition}{euristica-ancoraggio-definition}{Euristica dell'ancoraggio}
    L'\emph{euristica dell'ancoraggio} è una regola per cui i giudizi sul valore
    di un evento o di un esito mostrano aggiustamenti insufficienti, verso
    l'alto o verso il basso, rispetto a un valore di partenza.
\end{snippetdefinition}

\begin{snippetexample}{ancoraggio-prodotto-numeri-example}{Prodotto di numeri}
    Tversky e Kahneman chiesero ai partecipanti di stimare rapidamente il
    risultato di \(1 \times 2 \times 3 \times 4 \times 5 \times 6 \times 7
    \times 8\) oppure di \(8 \times 7 \times 6 \times 5 \times 4 \times 3
    \times 2 \times 1\). Le due operazioni hanno lo stesso risultato, ma le
    stime medie erano molto diverse perché i primi valori della sequenza
    funzionavano da ancore differenti.
\end{snippetexample}

\includesnpt[width=50\%,src=/snippet/static-psychology/anchoring-heuristic-example.jpg]{centered-img}

\begin{snippetdefinition}{punti-riferimento-definition}{Punti di riferimento}
    I \emph{punti di riferimento} sono criteri rispetto ai quali un esito viene
    interpretato come guadagno o come perdita. Dipendono in parte dalle
    aspettative del decisore.
\end{snippetdefinition}

\begin{snippetdefinition}{rammarico-definition}{Rammarico}
    Il \emph{rammarico} è l'emozione negativa che può seguire una decisione
    percepita come sbagliata. Quando è prevedibile, può rendere le persone più
    caute nella scelta.
\end{snippetdefinition}

\begin{snippet}{formulazione-decisione}
    Il modo in cui una domanda è formulata può influenzare le decisioni. Una
    stessa opzione può essere valutata diversamente se viene presentata rispetto
    a un guadagno, a una perdita o a un punto di riferimento diverso.
\end{snippet}

\end{document}
